\documentclass{article}
\usepackage[a4paper,margin=2.4cm]{geometry}
\usepackage{graphicx}
\usepackage{float}
\usepackage{fancyvrb}
\fvset{frame=single,fontsize=\small,framesep=3mm}

\title{Labwork 8: Scatter}
\author{Pham Tien Nhat}
\date{\today}

\begin{document}

\maketitle

\section{Implementation}

The input image of this labwork is stored as an Array of Struct (AoS) where each pixel keeps its R, G and B values together. In Labwork 8 I have converted the input image to HSV as a Struct of Array (SoA), with 3 different arrays $H[\,]$, $S[\,]$ and $V[\,]$, then I have also converted the HSV image version back to RGB original image.

\begin{Verbatim}
import math
import time
import numpy as np
import matplotlib.pyplot as plt

image = plt.imread('1.jpeg')
plt.imshow(image)
h, w, d = image.shape
\end{Verbatim}
\begin{figure}[H]
\centering
\includegraphics[width=0.6\textwidth]{../../image/lab8_input.png}
\caption{Input image.}
\label{figin}
\end{figure}

\textbf{RGB to HSV Conversion:} Each thread reads one pixel, scales R, G, B from $[0..255]$ to $[0..1]$, and finds $max$, $min$ and $\Delta = max - min$ following the formula given in the slide. The hue $H \in [0, 360]$ depends on which channel is the maximum, the saturation is $S = \Delta / max$ (0 if $max = 0$) and the value is $V = max$. The kernel is a scatter: one input pixel $(R, G, B)$ is written to 3 different output arrays.

\begin{Verbatim}
@cuda.jit
def RGB2HSV(src, H, S, V):
    tidx = cuda.threadIdx.x + cuda.blockIdx.x * cuda.blockDim.x
    tidy = cuda.threadIdx.y + cuda.blockIdx.y * cuda.blockDim.y
    if tidx >= src.shape[1] or tidy >= src.shape[0]:
        return

    r = src[tidy, tidx, 0] / 255
    g = src[tidy, tidx, 1] / 255
    b = src[tidy, tidx, 2] / 255
    mx = max(r, g, b)
    mn = min(r, g, b)
    delta = mx - mn

    # H conversion
    if delta == 0:
        hue = 0.0
    elif mx == r:
        hue = 60 * (((g - b) / delta) % 6)
    elif mx == g:
        hue = 60 * ((b - r) / delta + 2)
    else:
        hue = 60 * ((r - g) / delta + 4)

    # S conversion
    if mx == 0:
        sat = 0.0
    else:
        sat = delta / mx

    # scatter: one pixel (R, G, B) goes to 3 different arrays
    H[tidy, tidx] = hue
    S[tidy, tidx] = sat
    V[tidy, tidx] = mx

blockSize = (32, 32)
gridSize = (math.ceil(w / blockSize[0]), math.ceil(h / blockSize[1]))
devSrc = cuda.to_device(image)
devH = cuda.device_array((h, w), np.float32)
devS = cuda.device_array((h, w), np.float32)
devV = cuda.device_array((h, w), np.float32)
RGB2HSV[gridSize, blockSize](devSrc, devH, devS, devV)
\end{Verbatim}

The output of HSV process is display in the figure below, with three images each represents a channel from HSV, where $H$ represents the color changes, value is between 0 to 360, $S$ shows the colorfulness, value lies within 0 and 1 and $V$ displays the brightness, value ranging from 0 to 1

\begin{figure}[H]
\centering
\includegraphics[width=\textwidth]{../../image/lab8_hsv.png}
\caption{The 3 output arrays $H$, $S$ and $V$.}
\label{fighsv}
\end{figure}

\textbf{HSV to RGB Conversion:} The kernel reads the 3 arrays $H$, $S$, $V$ and writes one RGB image.

\begin{Verbatim}
@cuda.jit
def HSV2RGB(src, dst, H, S, V):
    tidx = cuda.threadIdx.x + cuda.blockIdx.x * cuda.blockDim.x
    tidy = cuda.threadIdx.y + cuda.blockIdx.y * cuda.blockDim.y
    if tidx >= src.shape[1] or tidy >= src.shape[0]:
        return

    # Preparation
    d = H[tidy, tidx] / 60
    hi = int(d) % 6
    f = d - hi
    l = V[tidy, tidx] * (1 - S[tidy, tidx])
    m = V[tidy, tidx] * (1 - f * S[tidy, tidx])
    n = V[tidy, tidx] * (1 - (1 - f) * S[tidy, tidx])

    # Conversion
    if 0 <= H[tidy, tidx] < 60:
        R = V[tidy, tidx]
        G = n
        B = l
    elif 60 <= H[tidy, tidx] < 120:
        R = m
        G = V[tidy, tidx]
        B = l
    elif 120 <= H[tidy, tidx] < 180:
        R = l
        G = V[tidy, tidx]
        B = n
    elif 180 <= H[tidy, tidx] < 240:
        R = l
        G = m
        B = V[tidy, tidx]
    elif 240 <= H[tidy, tidx] < 300:
        R = n
        G = l
        B = V[tidy, tidx]
    else:
        R = V[tidy, tidx]
        G = l
        B = m
    R = R * 255
    G = G * 255
    B = B * 255

    dst[tidy, tidx, 0] = R
    dst[tidy, tidx, 1] = G
    dst[tidy, tidx, 2] = B

devDst = cuda.device_array((h, w, 3), np.uint8)
HSV2RGB[gridSize, blockSize](devSrc, devDst, devH, devS, devV)
result = devDst.copy_to_host()
\end{Verbatim}


\section{Experimental Results}

Figure~\ref{figrgb} compares the input image with the output after RGB to HSV and RGB conversions.

\begin{figure}[H]
\centering
\includegraphics[width=0.9\textwidth]{../../image/lab8_rgb.png}
\caption{Input image (left) and output after RGB $\rightarrow$ HSV $\rightarrow$ RGB (right).}
\label{figrgb}
\end{figure}

The two images look similar to each other after the conversions step.

\end{document}
